% GENERATED FILE. Edit canonical lessons, then run npm run generate:books.
% canonical-source-hash: fnv1a64:ad1bba22a6e101a1
% canonical-lessons: FR-C503-jardinage, FR-C503-jardiner, FR-C503-semer, FR-C503-arroser, FR-C503-arrosoir

\chapter{Le jardinage}
\label{ch:fr-le-jardinage}
\hlchaptermodality{503}

\begin{chapteropening}
\textbf{By the end of this chapter:} \emph{I can talk about gardening.}

Complete the last lesson of chapter 503: I can talk about gardening.
\end{chapteropening}

\section[le jardinage]{le jardinage — gardening}

\label{lesson:FR-C503-jardinage}



\begin{quote}
Before the new one: say the French for to frame, then the French for tripod.
\end{quote}

\subsection*{You'll want to know: le jardinage}
\textbf{le jardinage} — \textquotedblleft{}gardening\textquotedblright{}.

\begin{quote}
Le jardinage me détend après le travail.
\end{quote}

\begin{quote}
\emph{Gardening relaxes me after work.}
\end{quote}

\emph{Faire du jardinage} — to do some gardening.

\begin{cousinweb}
\emph{le jardinage} is built on \emph{jardin} \textquotedblleft{}garden\textquotedblright{} with \emph{-age}; \emph{jardin} goes back to Frankish \emph{gardo} \textquotedblleft{}an enclosure\textquotedblright{}, a Germanic word.
\end{cousinweb}

\begin{grammarlens}[title={What you've built}]
One more word for telling, explaining and describing at B1.
\end{grammarlens}

\subsection*{Your turn}
\begin{itemize}
\raggedright
  \item \emph{Say it:} \emph{le jardinage}
  \item \emph{Say it:} \emph{le jardinage}, once more
  \item \emph{Say it:} a sentence of your own with \emph{le jardinage}
  \item \emph{Recall:} say the French for to frame, then the French for tripod, then say \emph{le jardinage} again
\end{itemize}

\begin{tcolorbox}[breakable,colback=teal!4,colframe=teal!35!black,title={Before you move on}]
What does le jardinage mean? (\textquotedblleft{}Gardening\textquotedblright{}.) Say it once more.
\end{tcolorbox}
\section[jardiner]{jardiner — to garden}

\label{lesson:FR-C503-jardiner}



\begin{quote}
Before the new one: say the French for tripod, then the French for gardening.
\end{quote}

\subsection*{You'll want to know: jardiner}
\textbf{jardiner} — \textquotedblleft{}to garden\textquotedblright{}.

\begin{quote}
Ma mère jardine tous les samedis.
\end{quote}

\begin{quote}
\emph{My mother gardens every Saturday.}
\end{quote}

A regular \emph{-er} verb.

\begin{cousinweb}
\emph{jardiner} is built on \emph{jardin} \textquotedblleft{}garden\textquotedblright{}, which goes back to Frankish \emph{gardo} \textquotedblleft{}an enclosure\textquotedblright{}.
\end{cousinweb}

\begin{grammarlens}[title={What you've built}]
One more word for telling, explaining and describing at B1.
\end{grammarlens}

\subsection*{Your turn}
\begin{itemize}
\raggedright
  \item \emph{Say it:} \emph{jardiner}
  \item \emph{Say it:} \emph{jardiner}, once more
  \item \emph{Say it:} a sentence of your own with \emph{jardiner}
  \item \emph{Recall:} say the French for tripod, then the French for gardening, then say \emph{jardiner} again
\end{itemize}

\begin{tcolorbox}[breakable,colback=teal!4,colframe=teal!35!black,title={Before you move on}]
What does jardiner mean? (\textquotedblleft{}To garden\textquotedblright{}.) Say it once more.
\end{tcolorbox}
\section[semer]{semer — to sow}

\label{lesson:FR-C503-semer}



\begin{quote}
Before the new one: say the French for gardening, then the French for to garden.
\end{quote}

\subsection*{You'll want to know: semer}
\textbf{semer} — \textquotedblleft{}to sow\textquotedblright{}.

\begin{quote}
En avril, nous semons des carottes.
\end{quote}

\begin{quote}
\emph{In April we sow carrots.}
\end{quote}

\emph{Je sème}, \emph{nous semons}. Informally, to lose someone: \emph{il nous a semés}.

\begin{cousinweb}
\emph{semer} is inherited from Latin \emph{sēmināre} \textquotedblleft{}to sow\textquotedblright{} (EtymDB), from \emph{sēmen} \textquotedblleft{}seed\textquotedblright{} (Lewis \& Short).
\end{cousinweb}

\begin{grammarlens}[title={What you've built}]
One more word for telling, explaining and describing at B1.
\end{grammarlens}

\subsection*{Your turn}
\begin{itemize}
\raggedright
  \item \emph{Say it:} \emph{semer}
  \item \emph{Say it:} \emph{semer}, once more
  \item \emph{Say it:} a sentence of your own with \emph{semer}
  \item \emph{Recall:} say the French for gardening, then the French for to garden, then say \emph{semer} again
\end{itemize}

\begin{tcolorbox}[breakable,colback=teal!4,colframe=teal!35!black,title={Before you move on}]
What does semer mean? (\textquotedblleft{}To sow\textquotedblright{}.) Say it once more.
\end{tcolorbox}
\section[arroser]{arroser — to water (plants)}

\label{lesson:FR-C503-arroser}



\begin{quote}
Before the new one: say the French for to garden, then the French for to sow.
\end{quote}

\subsection*{You'll want to know: arroser}
\textbf{arroser} — \textquotedblleft{}to water (plants)\textquotedblright{}.

\begin{quote}
Pendant les vacances, le voisin arrose nos plantes.
\end{quote}

\begin{quote}
\emph{During the holidays the neighbour waters our plants.}
\end{quote}

Informally, to celebrate with a drink: \emph{on va arroser ça}.

\begin{cousinweb}
\emph{arroser} comes from Late Latin \emph{arrōsāre} \textquotedblleft{}to bedew\textquotedblright{}, built on Latin \emph{rōs} \textquotedblleft{}dew\textquotedblright{} (Lewis \& Short).
\end{cousinweb}

\begin{grammarlens}[title={What you've built}]
One more word for telling, explaining and describing at B1.
\end{grammarlens}

\subsection*{Your turn}
\begin{itemize}
\raggedright
  \item \emph{Say it:} \emph{arroser}
  \item \emph{Say it:} \emph{arroser}, once more
  \item \emph{Say it:} a sentence of your own with \emph{arroser}
  \item \emph{Recall:} say the French for to garden, then the French for to sow, then say \emph{arroser} again
\end{itemize}

\begin{tcolorbox}[breakable,colback=teal!4,colframe=teal!35!black,title={Before you move on}]
What does arroser mean? (\textquotedblleft{}To water (plants)\textquotedblright{}.) Say it once more.
\end{tcolorbox}
\section[l'arrosoir]{l'arrosoir — watering can}

\label{lesson:FR-C503-arrosoir}



\begin{quote}
Before the new one: say the French for to sow, then the French for to water.
\end{quote}

\subsection*{You'll want to know: l'arrosoir}
\textbf{l'arrosoir} — \textquotedblleft{}watering can\textquotedblright{}.

\begin{quote}
L'enfant remplit l'arrosoir au robinet.
\end{quote}

\begin{quote}
\emph{The child fills the watering can at the tap.}
\end{quote}

Masculine.

\begin{cousinweb}
\emph{l'arrosoir} is built on \emph{arroser} \textquotedblleft{}to water\textquotedblright{} with \emph{-oir}, the ending for a tool; \emph{arroser} goes back to Latin \emph{rōs} \textquotedblleft{}dew\textquotedblright{} (Lewis \& Short).
\end{cousinweb}

\begin{grammarlens}[title={What you've built}]
One more word for telling, explaining and describing at B1.
\end{grammarlens}

\subsection*{Your turn}
\begin{itemize}
\raggedright
  \item \emph{Say it:} \emph{l'arrosoir}
  \item \emph{Say it:} \emph{l'arrosoir}, once more
  \item \emph{Say it:} a sentence of your own with \emph{l'arrosoir}
  \item \emph{Recall:} say the French for to sow, then the French for to water, then say \emph{l'arrosoir} again
\end{itemize}

\begin{tcolorbox}[breakable,colback=teal!4,colframe=teal!35!black,title={Before you move on}]
What does l'arrosoir mean? (\textquotedblleft{}Watering can\textquotedblright{}.) Say it once more.
\end{tcolorbox}
